% Figure 1-1 — Capsule history of types in computer science and logic
% Redrawn as a structured LaTeX table from the source figure.
% Italic entries are developments in logic, matching the original convention.
\begin{figure}[p]
  \centering
  \scriptsize
  \renewcommand{\arraystretch}{0.92}
  \begin{tabularx}{\textwidth}{
    @{}>{\raggedright\arraybackslash}p{1.25cm}
    >{\raggedright\arraybackslash}X
    >{\raggedleft\arraybackslash}p{6.1cm}@{}
  }
    \toprule
    年代 & 进展 & 参考文献 \\
    \midrule
    1870 年代 & \emph{形式逻辑的起源} & Frege（1879） \\
    1900 年代 & \emph{数学的形式化} & Whitehead 和 Russell（1910） \\
    1930 年代 & \emph{无类型 lambda 演算} & Church（1941） \\
    1940 年代 & \emph{简单类型 lambda 演算}
      & Church（1940）；Curry 和 Feys（1958） \\
    \addlinespace
    1950 年代 & Fortran & Backus（1981） \\
      & Algol-60 & Naur 等（1963） \\
    \addlinespace
    1960 年代 & \emph{Automath 项目} & de Bruijn（1980） \\
      & Simula & Birtwistle 等（1979） \\
      & \emph{Curry--Howard 对应} & Howard（1980） \\
      & Algol-68 & van Wijngaarden 等（1975） \\
    \addlinespace
    1970 年代 & Pascal & Wirth（1971） \\
      & \emph{Martin-Löf 类型论} & Martin-Löf（1973，1982） \\
      & \emph{System \(F\)、\(F^\omega\)} & Girard（1972） \\
      & 多态 lambda 演算 & Reynolds（1974） \\
      & CLU & Liskov 等（1981） \\
      & 多态类型推断 & Milner（1978）；Damas 和 Milner（1982） \\
      & ML & Gordon、Milner 和 Wadsworth（1979） \\
      & \emph{交类型} & Coppo 和 Dezani（1978） \\
      & & Coppo、Dezani 和 Sallé（1979）；Pottinger（1980） \\
    \addlinespace
    1980 年代 & NuPRL 项目 & Constable 等（1986） \\
      & 子类型化 & Reynolds（1980）；Cardelli（1984）；Mitchell（1984a） \\
      & 以存在类型表示抽象数据类型 & Mitchell 和 Plotkin（1988） \\
      & \emph{构造演算} & Coquand（1985）；Coquand 和 Huet（1988） \\
      & \emph{线性逻辑} & Girard（1987）；Girard 等（1989） \\
      & 有界量化 & Cardelli 和 Wegner（1985） \\
      & & Curien 和 Ghelli（1992）；Cardelli 等（1994） \\
      & \emph{爱丁堡逻辑框架} & Harper、Honsell 和 Plotkin（1992） \\
      & Forsythe & Reynolds（1988） \\
      & \emph{纯类型系统}
        & Terlouw（1989）；Berardi（1988）；Barendregt（1991） \\
      & 依赖类型与模块化 & Burstall 和 Lampson（1984）；MacQueen（1986） \\
      & Quest & Cardelli（1991） \\
      & \termfirst{效应系统}{effect system}
        & Gifford 等（1987）；Talpin 和 Jouvelot（1992） \\
      & \termfirst{行变量}{row variable}；可扩展记录
        & Wand（1987）；Rémy（1989） \\
      & & Cardelli 和 Mitchell（1991） \\
    \addlinespace
    1990 年代 & 高阶子类型化
      & Cardelli（1990）；Cardelli 和 Longo（1991） \\
      & 带类型的中间语言 & Tarditi、Morrisett 等（1996） \\
      & 对象演算 & Abadi 和 Cardelli（1996） \\
      & \termfirst{半透明类型}{translucent type}与模块化
        & Harper 和 Lillibridge（1994）；Leroy（1994） \\
      & 类型化汇编语言 & Morrisett 等（1998） \\
    \bottomrule
  \end{tabularx}
  \caption{计算机科学与逻辑学中的类型简史}
  \label{fig:chapter-01-type-history}
\end{figure}
